% GENERATED FILE. Edit canonical lessons, then run npm run generate:books.
% canonical-source-hash: fnv1a64:141eeb37cd6aadf8
% canonical-lessons: KA-C391-samanate, KA-C391-taratamya-madu, KA-C391-sosisu, KA-C391-sabalikarana, KA-C391-baddhate

\chapter{Fairness for All}
\label{ch:ka-fairness-for-all}
\hlchaptermodality{391}

\begin{chapteropening}
\textbf{By the end of this chapter:} \emph{I can say equality, to discriminate, to exploit, empowerment and commitment.}

Complete the last lesson of chapter 391: I can say equality, to discriminate, to exploit, empowerment and commitment.
\end{chapteropening}

\section[\texorpdfstring{\kn{ಸಮಾನತೆ} (samānate)}{ಸಮಾನತೆ (samānate)}]{\kn{ಸಮಾನತೆ} (samānate) — equality}

\label{lesson:KA-C391-samanate}



\begin{quote}
Before the new one: say the Kannada for official, then the Kannada for interference.
\end{quote}

\subsection*{You'll want to know: \kn{ಸಮಾನತೆ}}
\textbf{\kn{ಸಮಾನತೆ}} — \emph{samānate} — \textquotedblleft{}equality\textquotedblright{}.

\begin{quote}
\kn{ಗಂಡು} \kn{ಹೆಣ್ಣಿನ} \kn{ನಡುವೆ} \kn{ಸಮಾನತೆ} \kn{ಇರಬೇಕು}.
\end{quote}

\begin{quote}
\emph{gaṇḍu heṇṇina naḍuve samānate irabēku.}
\end{quote}

\begin{quote}
There must be equality between men and women.
\end{quote}

\kn{ಸಮಾನ}, equal, + -\kn{ತೆ}, -ness.

\begin{cousinweb}
\textbf{\kn{ಸಮಾನತೆ}} (\emph{samānate}) is borrowed from Sanskrit \emph{samānatā}, \textquotedblleft{}equality\textquotedblright{} (Monier-Williams), from \emph{samāna}, \textquotedblleft{}equal\textquotedblright{}.
\end{cousinweb}

\begin{grammarlens}[title={What you've built}]
One more word for giving reasons and opinions.
\end{grammarlens}

\subsection*{Your turn}
\begin{itemize}
\raggedright
  \item \emph{Say it:} \emph{samānate}
  \item \emph{Say it:} \emph{samānate}, once more
  \item \emph{Say it:} \emph{gaṇḍu heṇṇina naḍuve samānate irabēku.} — There must be equality between men and women.
  \item \emph{Recall:} say the Kannada for official, then the Kannada for interference, then say \emph{samānate} again
\end{itemize}

\begin{tcolorbox}[breakable,colback=teal!4,colframe=teal!35!black,title={Before you move on}]
What does \kn{ಸಮಾನತೆ} mean? (\textquotedblleft{}Equality\textquotedblright{}.) Say it once more.
\end{tcolorbox}
\section[\texorpdfstring{\kn{ತಾರತಮ್ಯ} \kn{ಮಾಡು} (tāratamya māḍu)}{ತಾರತಮ್ಯ ಮಾಡು (tāratamya māḍu)}]{\kn{ತಾರತಮ್ಯ} \kn{ಮಾಡು} (tāratamya māḍu) — to discriminate}

\label{lesson:KA-C391-taratamya-madu}



\begin{quote}
Before the new one: say the Kannada for interference, then the Kannada for equality.
\end{quote}

\subsection*{You'll want to know: \kn{ತಾರತಮ್ಯ} \kn{ಮಾಡು}}
\textbf{\kn{ತಾರತಮ್ಯ} \kn{ಮಾಡು}} — \emph{tāratamya māḍu} — \textquotedblleft{}to discriminate\textquotedblright{}.

\begin{quote}
\kn{ಆ} \kn{ಅಂಗಡಿಯವನು} \kn{ಗ್ರಾಹಕರಲ್ಲಿ} \kn{ತಾರತಮ್ಯ} \kn{ಮಾಡುವುದಿಲ್ಲ}.
\end{quote}

\begin{quote}
\emph{ā aṅgaḍiyavanu grāhakaralli tāratamya māḍuvudilla.}
\end{quote}

\begin{quote}
That shopkeeper does not discriminate among customers.
\end{quote}

Past (I …): \kn{ತಾರತಮ್ಯ} \kn{ಮಾಡಿದೆ}, \emph{tāratamya māḍide}.

\kn{ತಾರತಮ್ಯ} alone is discrimination.

\begin{cousinweb}
\textbf{\kn{ತಾರತಮ್ಯ} \kn{ಮಾಡು}} (\emph{tāratamya māḍu}) is built from Sanskrit \emph{tāratamya}, \textquotedblleft{}gradation, difference\textquotedblright{} (Monier-Williams), with \emph{māḍu}, \textquotedblleft{}to do\textquotedblright{}: to rank people as more and less.
\end{cousinweb}

\begin{grammarlens}[title={What you've built}]
One more word for giving reasons and opinions.
\end{grammarlens}

\subsection*{Your turn}
\begin{itemize}
\raggedright
  \item \emph{Say it:} \emph{tāratamya māḍu}
  \item \emph{Say it:} \emph{tāratamya māḍu}, once more
  \item \emph{Say it:} \emph{ā aṅgaḍiyavanu grāhakaralli tāratamya māḍuvudilla.} — That shopkeeper does not discriminate among customers.
  \item \emph{Recall:} say the Kannada for interference, then the Kannada for equality, then say \emph{tāratamya māḍu} again
\end{itemize}

\begin{tcolorbox}[breakable,colback=teal!4,colframe=teal!35!black,title={Before you move on}]
What does \kn{ತಾರತಮ್ಯ} \kn{ಮಾಡು} mean? (\textquotedblleft{}To discriminate\textquotedblright{}.) Say it once more.
\end{tcolorbox}
\section[\texorpdfstring{\kn{ಶೋಷಿಸು} (śōṣisu)}{ಶೋಷಿಸು (śōṣisu)}]{\kn{ಶೋಷಿಸು} (śōṣisu) — to exploit}

\label{lesson:KA-C391-sosisu}



\begin{quote}
Before the new one: say the Kannada for equality, then the Kannada for to discriminate.
\end{quote}

\subsection*{You'll want to know: \kn{ಶೋಷಿಸು}}
\textbf{\kn{ಶೋಷಿಸು}} — \emph{śōṣisu} — \textquotedblleft{}to exploit\textquotedblright{}.

\begin{quote}
\kn{ಆ} \kn{ಕಂಪನಿ} \kn{ಕೆಲಸಗಾರರನ್ನು} \kn{ಶೋಷಿಸಿತು}.
\end{quote}

\begin{quote}
\emph{ā kampani kelasagārarannu śōṣisitu.}
\end{quote}

\begin{quote}
That company exploited its workers.
\end{quote}

Past (I …): \kn{ಶೋಷಿಸಿದೆ}, \emph{śōṣiside}.

Who is exploited takes -\kn{ಅನ್ನು}. \kn{ಶೋಷಣೆ}, exploitation.

\begin{cousinweb}
\textbf{\kn{ಶೋಷಿಸು}} (\emph{śōṣisu}) is built on Sanskrit \emph{śoṣaṇa}, \textquotedblleft{}drying up, draining\textquotedblright{} (Monier-Williams), with the Kannada verb ending -\emph{isu}: to drain someone.
\end{cousinweb}

\begin{grammarlens}[title={What you've built}]
One more word for giving reasons and opinions.
\end{grammarlens}

\subsection*{Your turn}
\begin{itemize}
\raggedright
  \item \emph{Say it:} \emph{śōṣisu}
  \item \emph{Say it:} \emph{śōṣisu}, once more
  \item \emph{Say it:} \emph{ā kampani kelasagārarannu śōṣisitu.} — That company exploited its workers.
  \item \emph{Recall:} say the Kannada for equality, then the Kannada for to discriminate, then say \emph{śōṣisu} again
\end{itemize}

\begin{tcolorbox}[breakable,colback=teal!4,colframe=teal!35!black,title={Before you move on}]
What does \kn{ಶೋಷಿಸು} mean? (\textquotedblleft{}To exploit\textquotedblright{}.) Say it once more.
\end{tcolorbox}
\section[\texorpdfstring{\kn{ಸಬಲೀಕರಣ} (sabalīkaraṇa)}{ಸಬಲೀಕರಣ (sabalīkaraṇa)}]{\kn{ಸಬಲೀಕರಣ} (sabalīkaraṇa) — empowerment}

\label{lesson:KA-C391-sabalikarana}



\begin{quote}
Before the new one: say the Kannada for to discriminate, then the Kannada for to exploit.
\end{quote}

\subsection*{You'll want to know: \kn{ಸಬಲೀಕರಣ}}
\textbf{\kn{ಸಬಲೀಕರಣ}} — \emph{sabalīkaraṇa} — \textquotedblleft{}empowerment\textquotedblright{}.

\begin{quote}
\kn{ಮಹಿಳೆಯರ} \kn{ಸಬಲೀಕರಣಕ್ಕೆ} \kn{ಶಿಕ್ಷಣ} \kn{ಮುಖ್ಯ}.
\end{quote}

\begin{quote}
\emph{mahiḷeyara sabalīkaraṇakke śikṣaṇa mukhya.}
\end{quote}

\begin{quote}
Education is key to women's empowerment.
\end{quote}

\kn{ಸಬಲ} (strong) + -\kn{ಈಕರಣ} (making).

\begin{cousinweb}
\textbf{\kn{ಸಬಲೀಕರಣ}} (\emph{sabalīkaraṇa}) is a modern word built on Sanskrit \emph{sabala}, \textquotedblleft{}powerful, strong\textquotedblright{} (Monier-Williams), with -\emph{īkaraṇa}, \textquotedblleft{}making into\textquotedblright{}: making strong.
\end{cousinweb}

\begin{grammarlens}[title={What you've built}]
One more word for giving reasons and opinions.
\end{grammarlens}

\subsection*{Your turn}
\begin{itemize}
\raggedright
  \item \emph{Say it:} \emph{sabalīkaraṇa}
  \item \emph{Say it:} \emph{sabalīkaraṇa}, once more
  \item \emph{Say it:} \emph{mahiḷeyara sabalīkaraṇakke śikṣaṇa mukhya.} — Education is key to women's empowerment.
  \item \emph{Recall:} say the Kannada for to discriminate, then the Kannada for to exploit, then say \emph{sabalīkaraṇa} again
\end{itemize}

\begin{tcolorbox}[breakable,colback=teal!4,colframe=teal!35!black,title={Before you move on}]
What does \kn{ಸಬಲೀಕರಣ} mean? (\textquotedblleft{}Empowerment\textquotedblright{}.) Say it once more.
\end{tcolorbox}
\section[\texorpdfstring{\kn{ಬದ್ಧತೆ} (baddhate)}{ಬದ್ಧತೆ (baddhate)}]{\kn{ಬದ್ಧತೆ} (baddhate) — commitment}

\label{lesson:KA-C391-baddhate}



\begin{quote}
Before the new one: say the Kannada for to exploit, then the Kannada for empowerment.
\end{quote}

\subsection*{You'll want to know: \kn{ಬದ್ಧತೆ}}
\textbf{\kn{ಬದ್ಧತೆ}} — \emph{baddhate} — \textquotedblleft{}commitment\textquotedblright{}.

\begin{quote}
\kn{ಅವಳ} \kn{ಕೆಲಸದಲ್ಲಿ} \kn{ಬದ್ಧತೆ} \kn{ಕಾಣುತ್ತದೆ}.
\end{quote}

\begin{quote}
\emph{avaḷa kelasadalli baddhate kāṇuttade.}
\end{quote}

\begin{quote}
Commitment shows in her work.
\end{quote}

\kn{ಬದ್ಧ}, bound: \kn{ಮಾತಿಗೆ} \kn{ಬದ್ಧ}, true to one's word.

\begin{cousinweb}
\textbf{\kn{ಬದ್ಧತೆ}} (\emph{baddhate}) is built on Sanskrit \emph{baddha}, \textquotedblleft{}bound, tied\textquotedblright{} (CDIAL 9126), with -\emph{tā}, \textquotedblleft{}-ness\textquotedblright{}: being bound to something.
\end{cousinweb}

\begin{grammarlens}[title={What you've built}]
One more word for giving reasons and opinions.
\end{grammarlens}

\subsection*{Your turn}
\begin{itemize}
\raggedright
  \item \emph{Say it:} \emph{baddhate}
  \item \emph{Say it:} \emph{baddhate}, once more
  \item \emph{Say it:} \emph{avaḷa kelasadalli baddhate kāṇuttade.} — Commitment shows in her work.
  \item \emph{Recall:} say the Kannada for to exploit, then the Kannada for empowerment, then say \emph{baddhate} again
\end{itemize}

\begin{tcolorbox}[breakable,colback=teal!4,colframe=teal!35!black,title={Before you move on}]
What does \kn{ಬದ್ಧತೆ} mean? (\textquotedblleft{}Commitment\textquotedblright{}.) Say it once more.
\end{tcolorbox}
